
Spurious correlations degrade worst-group accuracy while maintaining high average accuracy. Standard empirical risk minimization on Waterbirds achieves 95\% average accuracy but below 60\% worst-group accuracy when waterbirds appear on water backgrounds 90\% of the time. Existing detection methods require knowing which spurious feature to search for. Gradient attribution approaches (Grad-CAM, Integrated Gradients) fail for unknown spurious correlations --- user studies demonstrate practitioners cannot detect spurious reliance using saliency maps alone.

Current mitigation methods fall into three categories. Annotation-based methods (GroupDRO, JTT) require ground-truth group labels, achieving 80-85\% worst-group accuracy on Waterbirds. Embedding-space methods (SCER) operate post-hoc via prototype refinement, reporting estimated 90\% worst-group accuracy (code unavailable for reproduction). Attribution-based detection methods (SPROD) explain which features drive predictions but cannot identify whether the prediction process exhibits abnormality without knowing what to look for.

This creates a detection-mitigation gap: no unified framework exists that both detects minority groups via gradient analysis and mitigates spurious reliance without annotations.

We propose extending GAIA (Gradient Abnormality for Out-of-Distribution Detection) from distribution shift to subpopulation shift. Minority samples create conflict when spurious shortcuts contradict core features. This conflict manifests in gradient computation process --- abnormality in how gradients scatter rather than what they point to. We validate methodology through synthetic experiments that prove pipeline correctness and proof-of-concept tests on toy datasets.

\textbf{Contributions:}

C1. Detection methodology validated via synthetic experiments: gradient abnormality pipeline (GradCAM $\rightarrow$ GAIA-Z $\rightarrow$ statistical testing) differentiates minority vs majority patterns when controlled synthetic gradients exhibit known differences (divergence 0.30, p<0.0001, Cohen's d=198.75). First application of GAIA to subpopulation shift. Real Waterbirds validation pending.

C2. Causal mechanism validation via synthetic tests: background augmentation reduces GAIA-Z by 39.4\% (p<0.001, d=2.96), supporting spurious-conflict hypothesis. Strong correlation (|$\rho$|=0.975, p<0.001) between GAIA divergence and worst-group accuracy in synthetic model sweep.

C3. Mitigation methodology via proof-of-concept: spatial gradient regularization framework improves worst-group accuracy by 23 percentage points on MNIST+Color (78\% vs 55\%, 1 seed, 2 epochs) without catastrophic accuracy drop. Full 5-seed experiment and Waterbirds validation deferred.

C4. Unified gradient-based framework: same abnormality mechanism enables both detection (identify minority samples) and mitigation (suppress spurious gradients during training), complementing embedding-based methods that operate post-hoc.

All validation results are from synthetic data or minimal proof-of-concept experiments. Real Waterbirds experiments require GPU compatibility resolution (PyTorch 2.1+ for H100 sm\_90 support) or CPU training (estimated 70-90 GPU hours). We position this as methodology contribution (Tier 3) with clear path to empirical validation (Tier 2).
